\section{Raw singleton epochs are a source closure problem}
\label{sec:raw-epoch-theorem}

The previous local analysis bounds one acyclic singleton batch polynomially
in its current grammar, and gives an exponential-in-depth bound for a succession
of batches. There is a stronger algorithmic alternative for an entire epoch
in which the only word transformations are raw singleton substitutions.
Such an epoch can be discovered on its immutable source incidence data and
compiled as one batch. The essential new statement is the equivalence with
\emph{every possible sequential raw elimination order}. Computing the least
closure of a family of implications is standard; matching acyclicity and
unique perfect matchings are standard as well~\cite{kozen,golumbic,dowling}.

\subsection{Precisely which epochs are covered}

Let $X$ be a finite set of positive generator labels, and let
$(R_1,\ldots,R_m)$ be a list of raw words in $X^{\pm1}$ represented by one
shared straight-line grammar. Empty and repeated slots are retained. Fix a
survivor set $S\subseteq X$, and put $E=X\setminus S$.

A raw singleton step chooses a previously unconsumed original slot $d$
whose current word is $U x^\epsilon V$, where $\epsilon\in\{1,-1\}$ and
the absolute letter $x$ occurs exactly once. It sets
\[
 x\longmapsto (VU)^{-\epsilon},
\]
substitutes this word and its formal inverse into every remaining slot,
and consumes slot $d$. Every substitution is performed on raw words:
there is no free reduction, cyclic reduction, proper-power root extraction,
Whitehead transformation, or relator multiplication between steps.
An acyclic raw batch can be linearized in dependency order, so is included.
The epoch succeeds for $S$ if it eliminates exactly $E$.

The frozen source itself may be the result of an already verified prefix
containing other moves. It is the transformations \emph{inside} the epoch
that are restricted. In particular, this theorem can be applied afresh at
each normalized presentation boundary. It does not combine the different
boundaries into one bounded-cost search.

For each source row $R_j$ and each generator $x$ occurring exactly once in
its raw spelling, form the implication
\[
 \operatorname{supp}(R_j)\setminus\{x\}\ \Longrightarrow\ x,
 \tag{H}\label{eq:raw-horn-rule}
\]
where support ignores signs but counts distinct labels. Let
$\operatorname{cl}(S)$ be the least set containing $S$ and closed under
these implications. Multiplicity one is essential: a zero exponent sum is
not absence, and a nonzero exponent sum of magnitude one does not imply a
unique raw occurrence.

\begin{theorem}[Complete source criterion for raw singleton reduction]
\label{thm:raw-closure-equivalence}
The following are equivalent:
\begin{enumerate}
\item Some raw singleton epoch reduces the frozen source to the survivor
      alphabet $S$.
\item Some acyclic singleton batch formed directly from original source
      slots eliminates $E$.
\item $\operatorname{cl}(S)=X$.
\end{enumerate}
If an epoch uses a specified donor-slot set $D$, with $|D|=|E|$, then it can
be flattened using exactly the same source slots $D$, after possibly
reassigning their eliminated generators. The flattened batch and that
epoch induce the same homomorphism $F(X)\to F(S)$ as free-group elements.
Their unreduced output strings need not agree.
\end{theorem}

The last assertion requires the same donor set. A generic closure search
over all source rows may select a different donor set from a previous
producer. Its resulting presentation is still Tietze equivalent to the
source, but its word images are not asserted equal to those of the previous
producer in $F(S)$. No equality of tied witnesses, subsequent search paths,
or finite-budget recognition outcomes is implied.

\subsection{The positive pivot lemma}

For a nonnegative integral square matrix $A$, write $\operatorname{per}(A)$
for its permanent, with permanent one for the empty matrix. Equivalently,
replace entry $a_{ij}$ by that many separately labelled parallel edges in a
bipartite graph; the permanent counts its edge-labelled perfect matchings.
The permanent is used only in the proof, not computed by the algorithm.

\begin{lemma}[Reversing a successful unit pivot]
\label{lem:positive-pivot}
Suppose $a_{dc}=1$. Delete row $d$ and column $c$ and set
\[
 b_{ij}=a_{ij}+a_{ic}a_{dj}\qquad(i\ne d,\ j\ne c).
\]
If $\operatorname{per}(B)=1$, then $\operatorname{per}(A)=1$.
\end{lemma}

\begin{proof}
Tag each edge counted by $b_{ij}$ either as an original direct edge from
$i$ to $j$, or as an ordered pair of original edges from $i$ to $c$ and
from $d$ to $j$. Consider the unique tagged perfect matching of $B$.
It contains at most one pair-created edge. Indeed, if it contained
pair-created edges $(i,j)$ and $(k,\ell)$, their source rows and destination
columns would be distinct. Exchanging the destination factors would give
pair-created edges $(i,\ell)$ and $(k,j)$ and a second tagged perfect
matching, a contradiction.

If it contains no pair-created edge, adjoin the unique edge $(d,c)$.
If it contains one, replace that edge by its two original component edges.
Either operation constructs a perfect matching of $A$.

Conversely, each perfect matching of $A$ maps injectively to a tagged
perfect matching of $B$. If it uses $(d,c)$, delete that edge. Otherwise
there is a unique row $i$ matched to $c$ and a unique column $j$ matched
from $d$; replace those two edges by the corresponding pair-created edge
$(i,j)$. All remaining edges are tagged direct. The number of pair-created
edges distinguishes the two cases, and their tags recover all original
occurrence choices. Thus $A$ has at most as many perfect matchings as
$B$. It has at least one by the preceding construction, proving the lemma.
\end{proof}

The converse of the lemma for an arbitrarily chosen pivot is false. For
example
\[
 A=\begin{pmatrix}1&1&1\\1&0&1\\1&0&0\end{pmatrix}
 \quad\hbox{has permanent }1,
\]
but pivoting its upper-left entry gives
$\left(\begin{smallmatrix}1&2\\1&1\end{smallmatrix}\right)$, whose permanent
is three. A poor raw elimination order can become stuck even when another
order would finish. The source closure algorithm avoids this order dependence.

\subsection{Flattening proof}

For a completed epoch, let $D$ be its consumed original slots and $E$ its
eliminated generators. Form the square source occurrence matrix
\[
 A_{jx}=\#\{\text{occurrences of }x\text{ or }x^{-1}\text{ in }R_j\},
 \qquad j\in D,\quad x\in E.
\]
At any stage restrict to the still-unconsumed rows of $D$ and still-live
columns of $E$. A raw singleton step has pivot one, and its new occurrence
matrix is exactly the positive pivot in Lemma~\ref{lem:positive-pivot}.
Deleting the unique pivot letter leaves $a_{dy}$ occurrences of every other
letter $y$ in its defining image. Each of the $a_{ix}$ occurrences in
another row contributes one such image. Inversion changes signs and order,
but not absolute occurrence counts. No cancellation subtracts occurrences.

The final matrix is empty. Applying the lemma backwards through every step
gives $\operatorname{per}(A)=1$. Let its unique perfect matching assign
source donor $R_{d(x)}$ to $x$. Every matched occurrence has multiplicity
one. Direct an edge $x\to y$ when $y\ne x$, $y\in E$, and $y$ appears in
$R_{d(x)}$. A directed cycle would allow cyclic exchange of the matched
columns and create a second perfect matching. Hence this definition graph
is acyclic. Solving the matched source equations in dependency order gives
a single acyclic batch eliminating $E$.

For the equality assertion, let $\phi:F(X)\to F(S)$ be the homomorphism
computed by the original epoch. A consumed equation becomes the identity
when it is eliminated, and stays the identity under subsequent
homomorphisms. Therefore $\phi(R_j)=1$ for every $j\in D$, and $\phi$ fixes
$S$. The acyclic matching gives a unique $F(S)$-valued solution to these
same equations with the values on $S$ fixed: determine each generator in
dependency order. Its homomorphism $\psi$ must consequently equal $\phi$.
In particular, freely reduced residual source words agree. This uniqueness
argument is in a free group, not in the monoid of unreduced signed strings.

For a concrete difference in raw spelling, take the two donor equations
\[
 R_1=xayb,\qquad R_2=xc,
\]
with $a,b,c$ surviving. Eliminating $x$ first from $R_1$, then $y$ from
the substituted $R_2$, sends $x$ to
$b^{-1}bc^{-1}aa^{-1}$ as a raw word. The source matching assigns $R_2$ to
$x$ and $R_1$ to $y$, and sends $x$ directly to $c^{-1}$. Both maps agree
in $F(a,b,c)$.

To finish Theorem~\ref{thm:raw-closure-equivalence}, an acyclic source
definition order is a sequence of productive applications of
\eqref{eq:raw-horn-rule}, so the closure contains every generator.
Conversely, record one source row whenever closure first adds a generator.
Its other support labels were already known, so the resulting definitions
are acyclic. A row cannot be used twice: after its first productive firing,
every member of its support is known. These distinct source rows therefore
give a valid acyclic Tietze batch. This proves all directions.

\subsection{Exact discovery without constructing intermediate words}

The implementation uses the maintained word arena's immutable presence and
repetition masks. At a concatenation $UV$ they satisfy
\[
 P_{UV}=P_U\cup P_V,\qquad
 Q_{UV}=Q_U\cup Q_V\cup(P_U\cap P_V).
\]
Consequently $P_R\setminus Q_R$ is precisely the set of raw singleton
generators. Two uses of the same child grammar still count twice. These
summaries never infer cancellation from exponent sums or from equality of
oppositely signed words.

From the summaries, build a sparse generator-to-source-slot incidence
index. For each row retain its number of currently unknown distinct
generators and the xor of their dense local indices. When the number is
one, the xor identifies the remaining generator. The row can fire exactly
when that generator is in its original singleton set. A queue processes
newly enabled rows. On learning a generator, decrement the unknown counts
and update the xors in its incident rows. A row can pass from two unknown
generators to one only once, so it enters the queue at most once; a queued
row can become stale by falling to zero before it is visited.

The invariant is that the counters and xors describe exactly the unknown
members of every row. Every firing is permitted by the implication rule.
When no row can fire, the known set is closed. By monotonicity it is the
least closed superset of the seeds. No word image, substring, inverse, or
new grammar node is created by this discovery pass.

There is an exact early failure condition for all singleton seeds. If
$r=|X|>1$ and no source row has support size at most two and at least one
raw singleton generator, then no singleton seed admits a first productive
firing. All singleton closures are trivial. This is a proof about this
algebraic query; it is not a nontrivial-knot test.

Let $I=\sum_j|\operatorname{supp}(R_j)|$. Once source metadata is available,
one closure uses $O(r+m+I)$ structural operations and $O(r+m+I)$ working
entries. The counters and local xors have $O(\log(r+1))$ bits. Trying all
$r$ one-element seeds therefore takes
\[
 O\bigl(r(r+m+I)\bigr)
\]
such operations, plus source summarization. This is expected dictionary
operation accounting for the Python implementation; label and identifier
bit costs must be charged separately. Persistent arena masks may use the
historical alphabet width rather than the current live rank. That existing
memory cost is not hidden in the sparse source index bound.

More generally, the fixed-$k$-survivor query can enumerate all
$\binom rk$ seed sets, for $0\le k\le r$, with total closure cost
\[
 O\!\left(\binom rk(r+m+I)\right).
\]
This also decides reducibility to \emph{at most} $k$ survivors: a smaller
successful seed set extends to size $k$, and closure is monotone. For
$k\ge r$ the answer is immediate by the empty epoch; $k=0$ is the single
empty-seed query. For $k<r$, absence of every source row with support at
most $k+1$ and a raw singleton similarly proves that no size-$k$ seed set
can make a first productive firing. Only the $k=1$ early guard is used in
the delivered planner.

On a freshly renumbered compact source of encoded length $n_0$, the
parameters $r,m,M$ are $O(n_0)$ and $I\le mr=O(n_0^2)$. The one-seed
enumeration thus uses $O(n_0^3)$ structural operations after preparation.
Computing presence/repetition bit vectors, enumerating their set bits,
and performing index arithmetic contribute polynomial bit costs. This
establishes an actual polynomial decision algorithm for the stated raw
query. It does not identify the current source size $n_0$ with the
crossing count of a diagram before an unbounded transformation prefix.

For polynomially encoded sources of size $n$ and $k=O(\log n)$, this is a
quasipolynomial algorithm for \emph{raw singleton reducibility to at most
$k$ survivors}. It supplies neither a general knot predicate on the
remaining $k$ generators nor a bound on the transformations needed to enter
this class. Only the one-seed producer is implemented in this delivery;
the general enumeration is a direct algorithmic corollary.

\subsection{Polynomial epoch compilation and exact endpoint handling}

Once a successful closure is known, emit the original slot/generator pairs
as one existing version-eight \texttt{elimination\_batch} move. The
maintained producer and independently implemented checker reconstruct its
word images and acyclic dependency graph from those source slots. The
checker does not trust the closure index, xor counters, seed-search result,
or permanent argument as certificate fields.

Let $M$ be the number of reachable source grammar nodes, $H\le M$ its
height, and $q\le r$ the number of eliminated generators. Two source
substring cuts per donor, shared inversion, and one topological circuit
compilation give the conservative grammar-node bound
\[
 O(M+qH+q)=O((r+1)(M+1)).
\]
The $m$ ordered root/slot handles and the encoded label and node-index
bits are charged separately; the display is not a claim that empty or
repeated source slots have no storage cost.
All lengths have polynomial bit size because this is a binary terminal
grammar of polynomial size. This conservative argument alone suffices for
polynomial compilation independent of the number of raw phases in the
original epoch. The previously proved disjoint-hole-path analysis can
strengthen the single-batch allocation bound, but no linear-time claim is
needed here. Donor occurrence scans and independent replay can revisit a
source grammar once per donor; linear output size would not make them
linear-time algorithms.

For the same donor set as a given epoch, one further polynomial compressed
free-reduction pass~\cite{schleimer} yields exactly the epoch's freely reduced semantic
output. The theorem does not retroactively bound all allocations made by
the old sequential implementation, and does not preserve its unreduced
intermediate roots. It supplies a polynomial replacement using the frozen
source and a single batch.

At rank one, every residual relator is a power of the survivor. The
existing endpoint evaluates all residual exponent sums and accepts only
when every one is zero. Discarding unselected rows would be unsound for
general presentations: for example $\langle a,b\mid ba^{-1},a^2\rangle$
has a successful one-seed closure but presents a finite cyclic group.
A knot verdict continues to require a source-bound presentation and the
unchanged independent endpoint check.

\subsection{Why the restriction is necessary}

Normalization genuinely enlarges the move class. Consider
\[
 \langle z,x,y\mid xy,\ xy^2\rangle,
 \qquad S=\{z\}.
\]
The two source words are freely reduced. Their eliminated-variable
occurrence matrix is
$\left(\begin{smallmatrix}1&1\\1&2\end{smallmatrix}\right)$, with
permanent three, and the source closure of $\{z\}$ is unchanged.
Eliminating $x$ from the first row nevertheless gives $x=y^{-1}$; the
second row becomes the raw word $y^{-1}yy$. Free reduction turns this into
$y$, enabling a second elimination. Thus an intervening cancellation
invalidates the positive-pivot invariant. The theorem correctly excludes
that step, even though the presentation is cyclic.

On untouched Wirtinger presentations, seed propagation overlaps the
established Wirtinger-number framework~\cite{wirtinger} and the repository's existing
two-meridian stage. Applying general raw support rules after a verified
presentation prefix is broader, but does not retain meridional provenance
of arbitrary current generators automatically. Known unknot diagrams with
large diagrammatic Wirtinger number~\cite{incompatibility} also rule out interpreting failure of
a bounded-seed source query as evidence of knottedness.

\subsection{Independent finite checks}

The standalone matrix audit exhausts all $3^9=19{,}683$ ternary $3\times3$
matrices and all $2^{16}=65{,}536$ binary $4\times4$ matrices, together with
the smaller ternary cases. It compares three independently expressed
conditions: permanent one by permutation enumeration, existence of a full
raw unit-pivot sequence by recursive search, and source-row peeling.
All agree. There are respectively 654 and 13,032 successful matrices in
the two largest families.

Two thousand random signed triangular source constructions yield 1,896
completed randomly ordered raw epochs. For each, independent literal
substitution and same-donor source flattening have identical freely reduced
generator images and residual words. The original-to-source donor matching
changes in 534 cases, and unreduced retained words differ in 419 cases.
These differences directly exercise the distinction in the theorem.

The maintained-arena module's ten focused tests compare 8,822 fixed-seed
closures on 400 random signed presentations with a separate literal
fixed-point oracle. On another 600 random presentations, exhaustive
one-seed search agrees on the first successful seed; 351 successful plans
pass both the maintained literal and compressed independent batch replayers,
giving 702 replay checks. Further tests cover repeated and empty source
slots, first-firing failure, cache source changes, strict field types,
resource interruption, external cancellation, and a power with 4,097-bit
expanded length while the expansion method is disabled. These are local
algebraic tests, not hard-knot timing measurements or a proof of general
quasipolynomial unknot recognition.
